\input{../tex-inputs/latex-leseansicht-vorspann}

\section[Arthur Schnitzler an Gustav Schwarzkopf, 1. 9. 1916]{L04400 Arthur Schnitzler an Gustav Schwarzkopf, 1. 9. 1916}
\nopagebreak\mylabel{L04400v}
\rehead{ }\normalsize\beginnumbering\briefempfaengerindex{Schwarzkopf, Gustav@\textsc{Schwarzkopf, Gustav}!zzzSchnitzler, Arthur@\emph{von Arthur Schnitzler}!1916-09-011@{1. 9. 1916}|(be}
\toendnotes[C]{\smallbreak\pagebreak[2]}
\correspDesc{Versand  durch Arthur Schnitzler am 1. 9. 1916 in Salzburg
\newline{}Übermittlung  am 1. 9. 1916 in Salzburg
\newline{}Erhalt  durch Gustav Schwarzkopf im Zeitraum [2. 9. 1916 – 6. 9. 1916?] in Wien}\toendnotes[C]{\smallbreak}
\Standort{DLA, A:Schnitzler, HS.1985.1.1897.}
\physDesc{Bildpostkarte, 573 Zeichen
\newline{}Handschrift: Bleistift, lateinische Kurrent
\newline{}Versand: Stempel: »\nobreak{}1. IX. 16\nobreak{}«.  }\toendnotes[C]{\smallbreak}\pstart{}{\pb}Hrn Gustav Schwarzkopf\pend{}\pstart{}Wien I\oindex{I., Innere Stadt@\textbf{I., Innere Stadt}|pw}\pend{}\pstart{}Tiefer Graben 17\oindex{Wien@\textbf{Wien}!I., Innere Stadt@\textbf{I., Innere Stadt}!Tiefer Graben 17@\textbf{Tiefer Graben 17}|pw}\pend{}{\bigskip}
\pstart
           \centering{}{\pb}\textcolor{gray}{\textbf{Salzburg\oindex{Salzburg@\textbf{Salzburg}|pw} vom Kapuzinerberg\oindex{Kapuzinerberg@\textbf{Kapuzinerberg}|pw}.}}\pend
           \vspace{1em}
\pstart
           {\pb}Salzburg\oindex{Salzburg@\textbf{Salzburg}|pw}{ }1. 9. 1916\pend
           \vspace{0.5em}
\pstart
           lieber Gustav, ich habe Olga\pwindex{Schnitzler, Olga 17.\,1.\,1882 Wien – 13.\,1.\,1970 Lugano@\textsc{Schnitzler, Olga} (17.\,1.\,1882 Wien – 13.\,1.\,1970 Lugano)|pw}{ }heute begleitet, sie fährt morgen früh zu Liesl\pwindex{Steinrück, Elisabeth 19.\,11.\,1885 – 7.\,4.\,1920 Partenkirchen@\textsc{Steinrück, Elisabeth} (19.\,11.\,1885 – 7.\,4.\,1920 Partenkirchen)|pw}{ }\introOben{}auf  8 Tage\introOben{}, ders glücklicherweise nicht so
               schlecht geht, wie man uns fürchten ließ: sie will Olga\pwindex{Schnitzler, Olga 17.\,1.\,1882 Wien – 13.\,1.\,1970 Lugano@\textsc{Schnitzler, Olga} (17.\,1.\,1882 Wien – 13.\,1.\,1970 Lugano)|pw} sogar nach München\oindex{München@\textbf{München}|pw}
               entgegenfahren.\pend
           
\pstart
           – Ich reise morgen nach Altaussee\oindex{Altaussee@\textbf{Altaussee}|pw} weg, wo wir
               noch bis Ende Sept. bleiben wollen. \label{K_L04400-1v}\edtext{Heini\pwindex{Schnitzler, Heinrich 9.\,8.\,1902 Hinterbrühl – 12.\,7.\,1982 Wien@\textsc{Schnitzler, Heinrich} (9.\,8.\,1902 Hinterbrühl – 12.\,7.\,1982 Wien)|pw} mit Fingi\pwindex{Jacobus, Elisabeth 4.\,8.\,1877 Trier – 10.\,1.\,1963 Wien@\textsc{Jacobus, Elisabeth} (4.\,8.\,1877 Trier – 10.\,1.\,1963 Wien)|pw} fährt}{\lemma{\textnormal{\emph{Heini mit Fingi fährt}}}\Cendnote{\textnormal{Vgl. A. S.: \emph{Tagebuch}, 14. 9. 1916.}}}\label{K_L04400-1} früher! – und noch hier
                  \introOben{}(in A.\oindex{Altaussee@\textbf{Altaussee}|pw})\introOben{} erhält Lili\pwindex{Cappellini, Lili 13.\,9.\,1909 Wien – 26.\,7.\,1928 Venedig@\textsc{Cappellini, Lili} (13.\,9.\,1909 Wien – 26.\,7.\,1928 Venedig)|pw} ein \label{K_L04400-2v}\edtext{funkelnagelneues Fräulein\pwindex{Pfister, Angela-Julianna 22.\,2.\,1885 Aletshausen – 19.\,6.\,1917 Probstdeuben@\textsc{Pfister, Angela-Julianna} (22.\,2.\,1885 Aletshausen – 19.\,6.\,1917 Probstdeuben)|pwv}}{\lemma{\textnormal{\emph{funkelnagelneues Fräulein}}}\Cendnote{\textnormal{Vgl. A. S.: \emph{Tagebuch}, 15. 9. 1916.}}}\label{K_L04400-2}. In A.\oindex{Altaussee@\textbf{Altaussee}|pw} alles
               beim Alten! viele Bekannte, – aber glücklicherweise auch genug Saiblinge\pend
           
\pstart
           \label{T_L04400-1v}\edtext{– Sein Sie vielmals gegrüßt. \textcolor{gray}{Grüssen Sie} Dr Max\pwindex{Schwarzkopf, Max 12.\,6.\,1857 Wien – 14.\,4.\,1928 ebd.@\textsc{Schwarzkopf, Max} (12.\,6.\,1857 Wien – 14.\,4.\,1928 ebd.)|pw}}{\lemma{\textnormal{\emph{– Sein … Max}}}\Cendnote{\textnormal{seitlich am rechten Textrand}}}\label{T_L04400-1}\pend
           
\pstart
           \label{T_L04400-2v}\edtext{Herzlichst Ihr{\\[\baselineskip]}\spacefill\mbox{Arthur}}{\lemma{\textnormal{\emph{Herzlichst IhrArthur}}}\Cendnote{\textnormal{oberhalb des Textes, verkehrt zum Text}}}\label{T_L04400-2}\pend
           \leftskip=0em{}\selectlanguage{ngerman}\endnumbering\briefempfaengerindex{Schwarzkopf, Gustav@\textsc{Schwarzkopf, Gustav}!zzzSchnitzler, Arthur@\emph{von Arthur Schnitzler}!1916-09-011@{1. 9. 1916}|)be}\mylabel{L04400h}
\begin{anhang}
\end{anhang}\newcommand{\dateiname}{L04400}\newcommand{\titel}{Arthur Schnitzler an Gustav Schwarzkopf, 1. 9. 1916}\newcommand{\editorInnen}{Herausgegeben von Jahnke, SelmaMüller, Martin Anton}\input{../tex-inputs/latex-leseansicht-abspann}
